\documentclass[11pt]{article}
\usepackage[margin=1.2in]{geometry}
\usepackage{amsmath,amssymb,amsthm}
\usepackage{xr}
\externaldocument{../paper2a_note/note}
\usepackage{hyperref}
\sloppy

\newtheorem{theorem}{Theorem}
\newtheorem{corollary}[theorem]{Corollary}
\newtheorem{lemma}[theorem]{Lemma}
\newtheorem{conjecture}[theorem]{Conjecture}
\newtheorem{proposition}[theorem]{Proposition}
\theoremstyle{remark}
\newtheorem{remark}[theorem]{Remark}

\newcommand{\E}{\mathbb{E}}
\newcommand{\Q}{\mathbb{Q}}

\title{Roots of $d$-matching polynomials: supplementary material\\
to the closed form and the refuted conjecture}
\author{Vico Bonfioli\thanks{Lean~4 sources: \texttt{RamaLean/} in the \texttt{rama-notes} repository. Correspondence: \texttt{vicobonfioli@gmail.com}.}}
\date{August 2026}

\begin{document}
\maketitle

\begin{abstract}
This note carries the material removed from \cite{PartI} when that paper was cut to its
spine: the two families in which the refuted conjecture is nonetheless provable, three
reductions, the biregular interval, the two-parameter family containing Hall’s
counterexample, the approximate-localization measurements, and the record of the searches
that failed to find the counterexample which eventually appeared. Nothing here is needed to
read \cite{PartI}. It is kept because the negative record is the part most easily lost, and
because a false conjecture that survived five independent searches is worth documenting.
\end{abstract}

\section{Two provable families, and three reductions}

\subsection{Subdivisions: the case of minimum degree two}

\begin{theorem}\label{thm:subdiv}
Let $H$ be $D$-regular, $D\ge2$, and let $S(H)$ be its subdivision, so that the
universal cover of $S(H)$ is the $(D,2)$-biregular tree with spectral gap
$\bigl(0,\sqrt{D-1}-1\bigr)$. Then every nonzero root of $\mu_{S(H)}$ has
absolute value at least $\sqrt{D-1}-1$; that is,
the Conjecture of \cite{PartI} holds for $S(H)$ at $d=1$. The constant is the spectral
gap of the universal cover; we do not claim it is attained, and by the strictness
of Heilmann--Lieb on finite graphs it is not.
\end{theorem}

Theorem~\ref{thm:subdiv} settles the case $\min(d,r)=2$ of Problem~1 of
\cite{SFM}.
Their uniformity condition allows $r\ge2$, and they observe that when $H$ is a
simple graph the incidence graph $B_H$ is exactly the subdivision $S(H)$; so
their question at $r=2$ asks precisely whether $\mu_{S(H)}$ has no root $\mu$
with $0<|\mu|<|\sqrt{D-1}-1|$. By their Theorem 1.8 this gives every
$D$-regular graph, viewed as a $(D,2)$-regular hypergraph, a left-sided
Ramanujan $2$-covering. The case $d=2$ follows as well, since $B_H$ is a
bipartite graph and $\mu_{B_H}$ does not distinguish its two sides: a
$(2,r)$-regular hypergraph has for incidence graph the subdivision of an
$r$-regular graph, and the quantity $|\sqrt{d-1}-\sqrt{r-1}|$ is symmetric in
$d$ and $r$. The case $\min(d,r)\ge3$ of their problem remains open in general, as does
the Conjecture of \cite{PartI} for $\min(a,b)\ge3$; Theorem~\ref{thm:kpq} below settles the
complete bipartite graphs inside that range, $K_{3,4}$ among them.

\begin{proof}
By the subdivision identity $\mu_{S(H)}(x)=x^{|E|-|V|}\,\mu_H(x^2-D)$
\cite{WWM}, the nonzero roots of $\mu_{S(H)}$ are the $x$ with $x^2=D+\theta$
for $\theta$ a root of $\mu_H$. Heilmann--Lieb gives $\theta\ge-2\sqrt{D-1}$,
so with $s=\sqrt{D-1}$ we get $x^2\ge s^2+1-2s=(s-1)^2$, whence
$|x|\ge s-1$. Sharpness: for cubic $H$ of large girth $\theta_{\min}\to-2\sqrt2$,
so the bound is approached.
\end{proof}

The inequality step is machine-checked in Lean~4 as \texttt{subdivision\_gap}.
The theorem is exactly Heilmann--Lieb for $H$ transported through the
subdivision identity, and the numbers match: for $H=K_4$, $\mu_{K_4}=x^4-6x^2+3$
gives smallest positive root $\sqrt{3-\sqrt{3+\sqrt6}\,}=0.8158$, against the
gap edge $0.4142$; the subdivisions of $K_{3,3}$, the prism, $Q_3$, the Wagner,
Petersen, Heawood and M\"obius--Kantor graphs give
$0.701,0.714,0.650,0.647,0.607,0.563,0.551$, decreasing toward the gap edge
with the girth and never reaching it.

For $K_{2,3}$ the computation was pushed to $d=6$ using
$\mu_{d,G}=\E[\chi_H]/\chi_G$ over $(d+1)$-covers, which is polynomial time per
cover, together with the reduction of the first free permutation to one
representative per cycle type; the ratios continue
$1.68$ at $d=5$ and no root in the gap at $d=6$. That last case is a caution
worth recording: in floating point it appeared to produce roots at
$\pm0.055$, well inside the gap, which an exact integer recomputation showed to
be an artefact of evaluating a degree-$30$ polynomial built from $35\times35$
characteristic polynomials in double precision.

 The last column decreases monotonically in $d$ in
every case, which is what the conjecture predicts: by
the closed form of \cite{PartI} in the $b_1=1$ case the roots equidistribute on
$\operatorname{spec}(T)$, so they should approach the gap edge without crossing
it. The evidence is therefore not merely that the bound holds but that it is
approached, though no computation here reaches the regime where the ratio is
close to $1$.

\subsection{Complete bipartite graphs}\label{sec:kpq}

The smallest graph in the open range $\min(a,b)\ge3$ is $K_{3,4}$, which is
$(4,3)$-biregular. The complete bipartite family is settled outright, by an exact
computation rather than by transport from a regular graph.

\begin{theorem}\label{thm:kpq}
Let $2\le p\le q$. Every root of $\mu_{K_{p,q}}$ lies in
\[
\{0\}\cup\pm\bigl[\sqrt{q-1}-\sqrt{p-1},\ \sqrt{q-1}+\sqrt{p-1}\bigr],
\]
the spectrum of the $(q,p)$-biregular tree. So the Conjecture of \cite{PartI} holds for
every $K_{p,q}$ at $d=1$.
\end{theorem}

\begin{proof}
Since $m_k(K_{p,q})=\binom pk\binom qk k!$, comparison of coefficients with
$L_p^{(\alpha)}(t)=\sum_k(-1)^k\binom{p+\alpha}{p-k}t^k/k!$ at $\alpha=q-p$ gives
\[
\mu_{K_{p,q}}(x)=x^{\,q-p}(-1)^p\,p!\;L_p^{(q-p)}(x^2),
\]
so the nonzero roots are the square roots of the zeros of a Laguerre polynomial.
Those zeros are the eigenvalues of the $p\times p$ symmetric tridiagonal Jacobi
matrix $J$ with $J_{kk}=2k+\alpha+1$ and $J_{k,k+1}=J_{k+1,k}=\sqrt{(k+1)(k+\alpha+1)}$,
$0\le k\le p-1$.

Put $c=p+q-2$ and $g(u)=\sqrt{u(u+\alpha)}-u$, and apply Gershgorin to $J-cI$. The
shifted diagonal in row $k$ is $2k-2p+3$, which is at most $-1$ for $k\le p-2$, so
the radius of an interior row is
\[
(2p-3-2k)+\sqrt{k(k+\alpha)}+\sqrt{(k+1)(k+1+\alpha)}\;=\;2p-2+g(k)+g(k+1),
\]
the index cancelling exactly. The step inequality
\begin{equation}\label{eq:amgm}
\sqrt{u(u+\alpha)}+1\;\le\;\sqrt{(u+1)(u+1+\alpha)}
\end{equation}
reduces, on squaring, to $\sqrt{u(u+\alpha)}\le u+\alpha/2$, which is
$\alpha^2\ge0$. It makes $g$ nondecreasing along the integers, so the radius is
largest at $k=p-2$, where it is $1+\sqrt{(p-2)(p-2+\alpha)}+\sqrt{(p-1)(q-1)}$ and
\eqref{eq:amgm} at $u=p-2$ bounds it by $2\sqrt{(p-1)(q-1)}$. The last row has
shifted diagonal $1$ and the single off-diagonal entry $\sqrt{(p-1)(q-1)}$, so its
radius lies within the same bound as soon as $(p-1)(q-1)\ge1$, which holds for
$p\ge2$. Hence every eigenvalue $t$ of $J$ satisfies
$|t-(p+q-2)|\le2\sqrt{(p-1)(q-1)}$; since
$\bigl(\sqrt{q-1}\pm\sqrt{p-1}\bigr)^2=(p+q-2)\pm2\sqrt{(p-1)(q-1)}$, that disc is
exactly the stated band for $x^2=t$. The factor $x^{q-p}$ contributes the root $0$,
which lies in the tree spectrum.
\end{proof}

The bound is sharp: at $p=q$ every interior Gershgorin row attains it, and the
conclusion is then the Heilmann--Lieb interval for the $p$-regular $K_{p,p}$. It
does need $p\ge2$: at $p=1$ the only zero is $q$ while the band degenerates to
$\{q-1\}$, the $(q,1)$-``biregular tree'' being the finite star $K_{1,q}$. For
$K_{3,4}$ the Laguerre factor is $t^3-12t^2+36t-24$, with zeros
$0.935,\ 3.30,\ 7.76$ against the band $[0.101,\ 9.899]$.

Inequality \eqref{eq:amgm}, the two row bounds, the endpoint identity and the
Gershgorin step are machine-checked in Lean~4 as \texttt{LaguerreBand.sqrt\_step},
\texttt{rowRadius\_le}, \texttt{lastRow\_le}, \texttt{band\_endpoints} and
\texttt{laguerre\_band}. The coefficient identification and the passage to the Jacobi
matrix are classical and remain by hand.

\subsection{Three further reductions}\label{sec:further}

The following arose from a conversation Chris Hall conducted with a language model and
shared with us in August 2026. We have checked them, and record them here with their
status, which differs sharply between the three. None is used elsewhere in this paper.

\paragraph{Localization in the maximal abelian cover.}
Put a circle variable on each of the $b_1(G)$ cotree edges and let $A_G(z)$ be the
resulting magnetic adjacency matrix. Haar integration over $\mathbb T^{b_1}$ annihilates
every permutation cycle of length at least three, since a simple cycle must use a cotree
edge and therefore carries a nontrivial character; fixed points and transpositions survive,
and those are exactly the matchings. Hence
\[
\mu_G(x)=\int_{\mathbb T^{b_1}}\det\bigl(xI-A_G(z)\bigr)\,dz .
\]
For $|z|=1$ the matrix $A_G(z)$ is Hermitian, so the integrand is real for real $x$; and a
continuous real function on a connected space whose average vanishes must vanish somewhere.
Therefore every matching root is an eigenvalue of some $A_G(z)$, that is
$\operatorname{Zeros}(\mu_G)\subseteq\operatorname{spec}(G^{\mathrm{ab}})$, and
the Conjecture of \cite{PartI} at $d=1$ becomes
\[
\operatorname{Zeros}(\mu_G)\cap\bigl(\operatorname{spec}(G^{\mathrm{ab}})\setminus
\operatorname{spec}(T)\bigr)=\emptyset .
\]
The realness of the integrand, the vanishing-average argument and the reduction are
machine-checked (\texttt{AbelianCover.det\_sub\_smul\_one\_im},
\texttt{exists\_zero\_of\_integral\_zero}, \texttt{root\_mem\_of\_average},
\texttt{reduction}). The cancellation over cycle types, left by hand above, reduces to one
statement: a sum of nonzero vectors with pairwise disjoint supports is nonzero. The cycles of
a permutation are vertex-disjoint, hence use disjoint edge sets, hence have homology classes
of disjoint support, and each class is nonzero because a cycle must use a cotree edge. That
lemma is \texttt{TorusGodsilGutman.sum\_ne\_zero\_of\_disjoint\_support}, and the sign and
averaging steps are re-proved in the same file in the compact form
\texttt{sign\_const\_of\_nonvanishing}, \texttt{nonvanishing\_average\_ne\_zero},
\texttt{mu\_ne\_zero\_of\_nonvanishing}, the last being the contrapositive actually used.
The two graph-theoretic inputs it consumes are also machine-checked. That every cycle uses at
least one cotree edge is \texttt{CotreeCycle.exists\_cotree\_edge}: if all its edges lay in
the tree, the walk would transfer to a cycle there, which acyclicity forbids, and no relation
between the tree and $G$ is needed beyond acyclicity. That vertex-disjoint cycles use disjoint
edge sets is \texttt{disjoint\_edges\_of\_disjoint\_support}. What remains informal is only
the identification of the monomial attached to a permutation with the sum of the homology
classes of its cycles, which is bookkeeping about how the magnetic matrix is set up rather
than a mathematical step. The
identity is checked numerically to $7\cdot10^{-15}$ on seven graphs with $b_1$ up to four
(\texttt{code/torus\_gg.py}).

It is worth being precise about what the connectedness of $\mathbb T^{b_1}$ contributes,
since Godsil--Gutman average over the $2^{|E|}$ signings and obtain the same identity. A
nonvanishing continuous function on a finite discrete set may take both signs, so no
conclusion about the average follows; on a connected space it may not. The inclusion
$\operatorname{Zeros}(\mu_G)\subseteq\operatorname{spec}(G^{\mathrm{ab}})$ is a consequence of
the topology of the parameter space, not of the averaging.

The localization is useful in the negative direction too: any counterexample must live in
$\operatorname{spec}(G^{\mathrm{ab}})\setminus\operatorname{spec}(T)$.

Two 2025 papers settle the companion question for the \emph{point} spectrum, in the opposite
direction. Li, Magee, Sabri and Thomas \cite{LMST} characterise the eigenvalues of
$G^{\mathrm{ab}}$ as the reals that are roots of $\mu_{G\setminus\Gamma}$ for every
$2$-regular subgraph $\Gamma$, which at $\Gamma=\emptyset$ gives
$\operatorname{pointspec}(G^{\mathrm{ab}})\subseteq\operatorname{Zeros}(\mu_G)$; Spier
\cite{Sp} proves that $G^{\mathrm{uni}}$ and $G^{\mathrm{ab}}$ have the same eigenvalues.
Those concern flat bands, so neither contains nor is contained in the inclusion above, which
concerns the full spectrum. But Spier's theorem says the two covers are indistinguishable at
the level of point spectrum, and the residue
$\operatorname{spec}(G^{\mathrm{ab}})\setminus\operatorname{spec}(T)$ is exactly where they
become distinguishable, so that is where a proof of the Conjecture of \cite{PartI} has to act.

\paragraph{Feedback vertex number one.}
Suppose some vertex $v$ has $F=G-v$ a forest, so that every cycle of $G$ passes through
$v$; the first Betti number is then unrestricted, and flowers, theta graphs and every
$K_{2,q}$ qualify. Writing the universal cover's adjacency operator as a matrix over
$C^*_r(F_{b_1})$ and eliminating the preimage of $F$, the Schur complement $S_v(x)$ onto
the fibre over $v$ is a \emph{finite} sum of group elements, so it lies in the group
algebra inside $C^*_r(F_{b_1})$ rather than merely in the von Neumann algebra
$L(F_{b_1})$. That is the load-bearing point, since $L(F_{b_1})$ is a $\mathrm{II}_1$
factor full of projections while $C^*_r(F_{b_1})$ is projectionless by Pimsner--Voiculescu.
A self-adjoint invertible element of a projectionless $C^*$-algebra is definite, so its
canonical trace cannot vanish. And that trace is
\[
\tau\bigl(S_v(x)\bigr)=x-\sum_{p\sim v}\frac{\mu_{F-p}(x)}{\mu_F(x)}
=\frac{\mu_G(x)}{\mu_{G-v}(x)} ,
\]
the first equality because returning to the same lift of $v$ forces a walk to enter and
leave a lifted component of $F$ at the same attachment point (two distinct attachment
points of one component sit under distinct lifts of $v$, else the universal cover would
contain a cycle), and the second by the matching recursion at $v$ together with the
cofactor formula for a forest resolvent. A root of $\mu_G$ that is not a root of
$\mu_{G-v}$ therefore makes $\tau(S_v)$ vanish, forcing $S_v$ to be non-invertible, hence
the root lies in $\operatorname{spec}(T)$.

We machine-check the pieces that can be: the combinatorial core
(\texttt{FeedbackVertex.no\_bypass}), the resolvent diagonal
(\texttt{resolvent\_diag}), the algebraic collapse of the trace formula
(\texttt{matching\_ratio}), and the trace step (\texttt{trace\_ne\_zero\_of\_definite}).
Pimsner--Voiculescu is not in Mathlib, so \texttt{feedback\_one\_of} carries it as an
explicit hypothesis in the shape the argument consumes. We therefore state this as
established modulo that classical input, and not as machine-checked end to end.

\paragraph{Two closure operations.}
Attaching a fixed rooted tree $R$ at every vertex sends $\mu_{d,H}$ to
$\beta^{N}\mu_{d,H}(\alpha/\beta)$ with $\alpha=\mu_R$, $\beta=\mu_{R-o}$ and
$N=d\,|V(H)|$, because an uncovered vertex leaves its whole copy of $R$ free while a
covered one leaves $R-o$; the same ratio $\alpha/\beta$ is what the Schur complement of
the universal cover produces, so the conjecture is inherited. Starting from a regular
graph, where the conjecture is immediate, this yields graphs of arbitrarily large first
Betti number whose universal covers have genuine gaps. Independently, the subdivision
identity survives averaging over $d$-covers, since suppressing the degree-two fibres of a
cover of $S(H)$ gives a cover of $H$ and a product of independent uniform permutations is
uniform; that upgrades Theorem~\ref{thm:subdiv} from $d=1$ to every $d$, and with it the
whole $(D,2)$-biregular family. The substitution identity and the root transfer are
machine-checked (\texttt{TreeSubstitution.homogenize}, \texttt{star\_ratio},
\texttt{subdivision\_root}, \texttt{subdivision\_abs}); the covering combinatorics, which
is the same fibre bookkeeping as \texttt{CoverCounts}, is by hand. Neither operation
touches the irreducible core: both reduce the spectral problem to a \emph{scalar} transfer
function, and the difficulty begins where elimination leaves a matrix-valued one.

\subsection{The biregular case is one interval, and we have one end of it}
Let $G$ be $(a,b)$-biregular bipartite, $s=\sqrt{a-1}$, $t=\sqrt{b-1}$. Since
$G$ is bipartite, $\mu_G(x)=x^{\varepsilon}f(x^2)$ for a polynomial $f$ and
$\varepsilon\in\{0,1,2,\dots\}$. Put
\[
c=(a-1)+(b-1),\qquad m=(a-1)(b-1),\qquad g(z)=f(z+c).
\]
The universal cover has spectrum $\{0\}\cup\pm[\,|s-t|,s+t\,]$, and
$(s\pm t)^2=c\pm2\sqrt m$, so
\[
x\in\operatorname{spec}(T)\ \setminus\{0\}
\iff
x^2-c\in\bigl[-2\sqrt m,\ 2\sqrt m\bigr].
\]
The change of variable therefore turns the Conjecture of \cite{PartI}, for biregular
$G$, into a single clean statement:
\begin{equation}\label{eq:twosided}
\text{every root of }g\text{ lies in }\bigl[-2\sqrt m,\,2\sqrt m\bigr],
\end{equation}
which is a Heilmann--Lieb interval of effective degree $m+1$.

It is worth recording what this says about the matching generating polynomial
$M_G(w)=\sum_km_kw^k$, since that, rather than $\mu_G$, is the object of
\cite{HL}. From $f_G(y)=y^pM_G(-1/y)$ the roots correspond by $w=-1/y$, so a
lower bound on the least positive root of $f_G$ is an upper bound on the modulus
of the roots of $M_G$. Writing $Q_G(z)=M_G(-z)=\sum_k(-1)^km_kz^k$, the
reciprocal-and-reflect of $f_G$:
\begin{equation}\label{eq:annulus}
\text{the Conjecture of \cite{PartI} at }d=1
\iff
\text{every root of }Q_G\text{ lies in }\Bigl(0,\ (s-t)^{-2}\Bigr].
\end{equation}
Heilmann--Lieb, in the same variable, says every root of $Q_G$ is at least
$\bigl(2\sqrt{\Delta-1}\bigr)^{-2}$ in modulus. The two statements bound the
roots of one polynomial from the two sides, so \eqref{eq:twosided} and
\eqref{eq:annulus} are the same dichotomy seen in reciprocal coordinates. This
is a restatement and nothing more, but it is the coordinate in which the
classical machinery is written, and $Q_G$ is the polynomial whose deletion
recursion $Q_G=Q_{G-v}-z\sum_{u\sim v}Q_{G-v-u}$ is formalized in the
accompanying Lean development. The upper end of
\eqref{eq:twosided} is known: the path-tree argument of Godsil bounds the roots
of $\mu_G$ by the spectral radius of the path tree, which for an
$(a,b)$-biregular graph with $a,b\ge2$ is at most $s+t$
\cite{God81,God93}. It is worth
noting that \cite{HL} proved the corresponding statement only in the regular
case; the irregular case is due to Godsil and to \cite{MSS}. The lower end is
the conjecture. \emph{They are the two ends of one interval}, so
what is missing is exactly half of a two-sided Heilmann--Lieb theorem for
biregular bipartite graphs.



\section{Hall’s counterexample: the family, and how close the misses come}

\subsection{A family, and the role of the leaves}\label{sec:hallfamily}

Hall's graph is the case $p=q=5$ of the family $G(p,q)$ with $p$ branches, each a copy of
$K_{2,q}$ carrying a pendant leaf. The branch arithmetic is uniform: the last factor of
$\mu_{G(p,q)}$ is
\[
   x^{q}\bigl(x^4-(2q+1+p)x^2+q^2+p(q+1)\bigr).
\]
Whether its smaller root lies in a gap of $\operatorname{spec}(T)$ is a separate question
from whether that root is rational, and the two are independent: over $2\le p,q\le7$ the
discriminant is a perfect square in six cases, while nine cases give a root inside a gap,
and the two sets overlap in only three. The counterexamples are
\[
   (p,q)\in\{(7,3),(6,4),(7,4),(5,5),(6,5),(7,5),(6,6),(7,6),(7,7)\},
\]
with $41,43,43,49,50,55,57,64,71$ vertices; Hall's $(5,5)$ is the smallest. The least $p$
that works is $7,6,5,6,7$ for $q=3,4,5,6,7$, so the phenomenon is cheapest exactly at
$p=q=5$.

Removing the pendant leaves altogether, so that every $w_i$ has degree $q$, destroys the
effect over the same range: for all $36$ pairs with $2\le p,q\le7$ the root lies inside
$\operatorname{spec}(T)$, with the diagnostic separated by three orders of magnitude
(\texttt{code/hall\_family.py}). The test throughout is the scaling of the density of states
at the root, which vanishes linearly in $\eta$ outside the spectrum and tends to a positive
limit inside, and which avoids the cavity mass gate that these flat-band graphs fail.

That observation is real but it does not mean what it appears to mean. Neither minimum
degree at least two nor bounded maximum degree repairs the Conjecture of \cite{PartI}, and a
single graph on $31$ vertices, smaller than Hall's, refutes both at once.

\begin{proposition}\label{prop:subcubic}
There is a simple connected graph on $31$ vertices, with every degree equal to two or three,
having a root of its matching polynomial in an internal gap of $\operatorname{spec}(T)$.
\end{proposition}

Take the rooted binary tree of depth three and identify each of its eight leaves with a
vertex of an attached triangle. Internal skeleton vertices have degree three, each attachment
vertex has degree three, the remaining triangle vertices have degree two, and the skeleton
root has degree two, so the degree set is exactly $\{2,3\}$ and there are $31$ vertices and
$38$ edges. The matching polynomial has the factor $x^4-7x^2+8$, whose smaller positive root
\[
   \theta=\sqrt{\tfrac{7-\sqrt{17}}{2}}=1.19935282014558573\ldots
\]
lies in an internal gap of $\operatorname{spec}(T)$: the gap is narrow, of width about
$0.016$, with spectrum on both sides. The factorisation was computed twice by independent
routes, the rooted pair recursion for a tree of blocks and a memoised vertex-deletion
recursion, and the spectral exclusion is certified as in \cite{PartI}: Newton at $40$
digits, seeded from the cavity fixed point, gives a real ratio system on all $76$ directed
edges with residual $5\cdot10^{-41}$ and every ratio nonzero, with decay rate
\[
   \rho=0.99273509760803033\ldots<1 .
\]
The margin $1-\rho=0.0073$ is thinner than elsewhere in this section but is resolved far
beyond the precision used. Verified in \texttt{code/subcubic.py}.

The polynomial half is machine-checked in \texttt{RamaLean/SubcubicCounterexample.lean}. The
quartic is not an accident of the depth: at depth two the recursion already factors as
\[
   X\!\cdot\!A_1-2B_1=X^2(X^2-3)(X^4-7X^2+8)
\]
(\texttt{quartic\_factor}), so the quartic divides $A_2$ and hence $\mu_G=A_3$
(\texttt{quartic\_dvd\_A3}), whose degree is confirmed to be $31$
(\texttt{degree\_A3}). What depth three supplies is not the root but the gap: at depth two the
same root lies inside $\operatorname{spec}(T)$, and only the deeper skeleton opens a gap
around it. That is worth stating because it separates the two halves of a counterexample
cleanly, and says which one is the binding constraint here.

The construction explains why the earlier families needed large degrees: there the number of
branches was the degree of the central vertex, so many branches forced high degree. A tree
skeleton separates the two, supplying arbitrarily many branches at degree three. What the
mechanism needs is branches, not degree.

\begin{proposition}\label{prop:mindeg2}
There are simple connected bipartite graphs of minimum degree two with a root of the matching
polynomial in an internal gap of $\operatorname{spec}(T)$.
\end{proposition}

Replace each pendant leaf of $G(p,q)$ by a pendant cycle $C_k$ attached at $w_i$. Every new
vertex then has degree two, as do the middle vertices of $K_{2,q}$, so the graph has minimum
degree two, and the branches still meet only at $c$, so the deletion recurrence still
computes $\mu_G$. The effect survives, at a price in size: it needs larger $p$ and $q$ than
the pendant-leaf version, which is why a search stopping at $p,q\le7$ sees it only for
$k=5$. Counterexamples occur at
\[
   (k,p,q)\in\{(5,7,7),(4,7,8),(3,8,8),(5,7,8),(5,8,7),(4,8,8),(5,8,8)\},
\]
the smallest having $92$ vertices, against $41$ for Hall's. So the leaves make the
construction cheaper without being necessary to it.

For $(k,p,q)=(5,7,7)$, on $92$ vertices with $140$ edges and degrees between two and nine,
the root is the algebraic number $\theta=2.852901286261237\ldots$, the second largest root of the
irreducible octic $x^8-26x^6+187x^4-350x^2+91$, which divides $x\mu_H-p\mu_{H-v}$ and hence
$\mu_G$.

The spectral exclusion is certified, in a different form from \cite{HallCE} and for a
structural reason. A ratio system exists for \emph{every} $\lambda$ outside
$\operatorname{spec}(A_T)$, so the ratio equations impose no condition on $\lambda$ and there
is nothing to eliminate; exact algebra in a fixed number field is therefore not available in
the way it is for Hall's example, where the ratios happen to lie in
$\mathbb{Q}(\sqrt5,\sqrt{41})$. Instead the certificate is a posteriori. The $280$ directed
edges fall into $11$ orbits, listed in \texttt{code/mindeg2\_exact.py} together with their
non-backtracking follower counts, and the orbit equations are checked against the full
$280$-edge fixed point. Newton at $80$ digits gives ratios with residual
$5.6\cdot10^{-81}$, all nonzero; since $\|J^{-1}\|_\infty=204.8$, the Newton--Kantorovich
bound places an exact solution within $1.1\cdot10^{-78}$ of the computed one. The decay rate
there is
\[
   \rho=0.9666941618858803717\ldots,\qquad 1-\rho=0.0333,
\]
while the entrywise Lipschitz constant of the orbit quotient in the ratios is $45$, so $\rho$
can move by at most $5.7\cdot10^{-76}$ over the certified ball: the margin exceeds that by a
factor $6\cdot10^{73}$. Hence $\rho<1$ for the exact solution, and by \cite{AFH}
$\theta\notin\operatorname{spec}(A_T)$. A Collatz--Wielandt vector gives the same conclusion
independently, at $0.9829$. The statement is thus rigorous modulo Newton--Kantorovich exactly
as it is rigorous modulo \cite{AFH}, both classical. Verified in \texttt{code/mindeg2.py} and
\texttt{code/mindeg2\_exact.py}.

\subsection{Approximate localization}\label{sec:defect}

The counterexamples miss by very little. Locating the gap edges sharply, by bisecting the
Angel--Friedman--Hoory decay rate to $\rho=1$ rather than by scanning a density of states,
gives for each of them the distance from the root to $\operatorname{spec}(T)$:

\begin{center}
\begin{tabular}{lrrr}
family & $n$ & root & defect\\\hline
pendant leaf $(7,3)$ & $43$ & $1.88040$ & $0.0072$\\
pendant leaf $(5,5)$, Hall's & $41$ & $2.23607$ & $0.0118$\\
pendant leaf $(7,6)$ & $64$ & $2.47528$ & $0.0348$\\
pendant $C_5$ $(7,7)$ & $92$ & $2.85290$ & $0.0014$\\
pendant $C_3$ $(8,8)$ & $97$ & $3.01489$ & $0.0021$\\
triangles, depth $3$ & $31$ & $1.19935$ & $0.00024$\\
\end{tabular}
\end{center}

Over thirteen counterexamples on $31$ to $97$ vertices, drawn from two unrelated families, the
defect never exceeds $0.035$ and shows no growth in $n$. Every violated gap has width at most
$0.22$. This suggests that the Conjecture of \cite{PartI} is wrong in its constant rather than in
its shape.

\begin{conjecture}[D1]\label{conj:defect}
There is an absolute constant $C$ such that for every finite graph $G$ and every root $\theta$
of $\mu_G$, $\operatorname{dist}(\theta,\operatorname{spec}(T_G))\le C$.
\end{conjecture}

The statement was fixed before the confirming data was produced and has survived an
out-of-sample test on ten counterexamples disjoint from the three that suggested it. The
scaled form with $C/|V(G)|$ is not refuted but is not supported either: $\operatorname{defect}
\cdot n$ varies by a factor of three hundred with no pattern.

What D1 would buy is machine-checked in \texttt{RamaLean/DefectLocalization.lean}. It says
exactly that the roots lie in the closed $C$-neighbourhood of $\operatorname{spec}(T)$
(\texttt{zeros\_subset\_thickening}), so at $C=0$ it is the Conjecture of \cite{PartI} itself
(\texttt{thickening\_zero}); and it restores that conjecture away from the edges:

\begin{proposition}\label{prop:core}
If every root lies within $C$ of $\operatorname{spec}(T)$ and $(a,b)$ is a gap, then no root
lies in $(a+C,\,b-C)$.
\end{proposition}

So under D1 every gap wider than $2C$ has a root-free core, and any counterexample is confined
to a narrow gap, which is what the measured widths show. The open interval is not an artefact:
a root may sit exactly at distance $C$ from the spectrum, and the gap edges themselves belong
to $\operatorname{spec}(T)$ (\texttt{root\_free\_core}, \texttt{conj10\_of\_zero\_defect}).

This is the form in which the question seems worth attacking. It is also the natural one: a
root descending to a spectral edge but never passing it is the Alon--Boppana picture, and
the Conjecture of \cite{PartI} was the assertion that matching roots are Ramanujan in that sense
without exception. D1 says the exceptions exist but are uniformly small.


\section{The evidence that made D3 look true}

The conjecture D3 of \cite{PartI}, that minimum degree three suffices, is false. What follows
is the evidence that made it look true, kept because five independent searches came back clean
and the reason they did is instructive.

The evidence that made Conjecture D3 of \cite{PartI} plausible is recorded below rather than
deleted, because a false conjecture that survives five independent searches is worth a note on
how it managed it. First, the regular case is trivial and worth separating
off: for $d$-regular $G$ the universal cover is the $d$-regular tree, whose spectrum is exactly
the interval $[-2\sqrt{d-1},2\sqrt{d-1}]$ in which Heilmann and Lieb confine every root, so there
is nothing to violate (\texttt{conj10\_of\_regular}). The content lies in irregular graphs.

Second, minimum degree three narrows the gaps but does not remove them, so
Conjecture D3 of \cite{PartI} is not true for want of anywhere to hide: over $38$ graphs the mean
total gap width falls from $0.99$ at $\delta\le2$ to $0.09$ at $\delta\ge3$, while individual
gaps at $\delta=3$ reach $2.9$ (\texttt{code/gapscale.py}).

Third, the evidence covers a family unlike the rest. Wheels, a hub joined to a cycle, have
minimum degree three, are non-bipartite by construction, and carry gaps from $(2.02,2.08)$ at
$W_5$ out to $(2.02,3.20)$ at $W_{11}$. None up to $W_{14}$ has a root in any gap, and at the
closest approach the verdict rests on the density-of-states ladder rather than on locating a gap
edge: for $W_{10}$ through $W_{20}$ the largest root below $2$ has density of states constant to
four significant figures across four decades of $\eta$, so it lies in the spectrum
(\texttt{code/wheels.py}).

Fourth, and this is what the evidence amounted to. Five searches were run against
Conjecture D3 of \cite{PartI}, each designed against the failure mode of the last, and all came
back clean; the $14$-vertex counterexample of \cite{PartI} shows they were all clean for one reason, that each
raised the block's degrees by adding edges inside it. Two reproduce Hall's own mechanism at
minimum degree three --- $p$ copies of a block glued
at a new centre, so that the branch factorization makes every root of $\mu_H$ a root of $\mu_G$
--- first with off-the-shelf blocks and then with blocks built to the shape of his, $K_{2,q}$ with
a graph on the $q$-side lifting those degrees to three. A third abandons the construction
entirely: random graphs of minimum degree three with deliberately spread degrees and \emph{no cut
vertex}, which is the case a cut-based engine structurally cannot reach. A fourth requires a
separating pair. A fifth holds the degree sequence fixed, so the gap does not move, and varies
only the block's internal structure, so the roots do.

\begin{center}
\begin{tabular}{lrl}
\emph{engine} & \emph{scale} & \emph{outcome}\\[2pt]
$1$-cut, off-the-shelf blocks      & $118$ & closest approach $0.169$\\
$1$-cut, blocks of Hall's shape    & $350$ & clean\\
separator-free random search       & $419$ & clean\\
separating-pair random search      & $39$  & clean\\
fixed degrees, varying structure   & $14$  & clean
\end{tabular}
\end{center}

About twelve hundred configurations, and the two closest approaches are both cut-based
(\texttt{code/D1cut.py}, \texttt{code/D1cut\_adv.py}, \texttt{code/D3broad.py},
\texttt{code/D3fixdeg.py}). The nearest miss is worth stating because it is nearer than the table
suggests: a block $K_{2,12}$ with a perfect matching on the $12$-side has
$\mu_H=(x-1)^4(x+1)^4(x^6-26x^4+157x^2-12)$, so $1$ is a root of multiplicity four and hence a
root of $\mu_G$, and $1$ lies well inside the central gap of half-width $\sqrt{11}-\sqrt2=1.90$.
It is not a counterexample only because $1$ is a spectral \emph{atom} of the cover, so it lies in
$\operatorname{spec}(T)$ after all: the density of states there reads $8.97$, $89.7$ and $897$ as
$\eta$ falls through $10^{-2}$, $10^{-3}$, $10^{-4}$, the signature of a point mass. Breaking the
symmetry that creates the atom moves the root off it --- and into the continuous spectrum, not
into the gap, since the same change moves the degrees that place the gap.

The separator-free class is not merely clean but comfortable, which is a different statement and
the one worth having. Taking the density of states at a root as the margin, since a violation is
exactly a zero there, the median over $419$ graphs is stable at $0.15$ from $n=10$ to $n=20$ and
if anything rises; the \emph{minimum} falls like $1/n$, but the number of roots sampled doubles
over that range and a minimum over more draws is smaller for that reason alone, the observed
$4.19\times10^{-2}$ against the $4.04\times10^{-2}$ that sample size predicts
(\texttt{code/D3margin.py}). So Conjecture D3 of \cite{PartI} is not surviving these tests by a
hair.

There is one reformulation, and it comes from the same measurements rather than from a new idea.
Godsil gives $\mu_G\mid\mu_P$ for the path tree $P$, and $P$ is a tree, where the matching
polynomial is the characteristic polynomial; so every root of $\mu_G$ is an eigenvalue of $P$, and
$\operatorname{spec}(P)\subseteq\operatorname{spec}(T)$ implies the conjecture
(\texttt{PathTreeRoute.conj\_of\_pathtree}). The hypothesis is strictly stronger --- $\mu_P$
carries its own roots too --- and it fails for Hall's graph exactly where the conjecture does. The
containment holds for the biregular families measured, but at minimum degree three it is
\emph{false}: one graph in $86$ carrying a gap has a path tree with twelve eigenvalues inside it
while no root of $\mu_G$ is there. So the hypothesis fails exactly where
Conjecture D3 of \cite{PartI} holds, and the reduction cannot carry a proof of it; the companion
paper \cite{PartII} gives the counterexample. We record it because it closes a route that the
biregular evidence makes look plausible.

What none of it supplies is an argument. Each engine was built against the previous failure mode
and each came back clean, which raises confidence and yields nothing proof-shaped. Every
counterexample known has a separation, and the evidence is strongest exactly where separations are
absent; turning that observation into a theorem is the open problem, and it is not one more search.

\subsection{A second conjecture falls with it}\label{sec:xu}

The counterexample refutes more than the conjecture it was built for. Zili Xu \cite{Xu}
proposes a spectral-inclusion conjecture for \emph{additive products}
in the sense of Mohanty--O'Donnell \cite{MO}: with $A_1,\dots,A_c$ adjacency matrices on a common
vertex set, $X=\operatorname{AddProd}(A_1,\dots,A_c)$, and
\[
   \alpha(A_1,\dots,A_c;x)=\mathbb E\det\Bigl(xI-\sum_{j}Q_j^{*}A_jQ_j\Bigr)
\]
the additive characteristic polynomial, the expectation over independent uniform $\pm1$ diagonal
signings, he conjectures $Z(\alpha)\subseteq\operatorname{Spec}(X)$. Mohanty--O'Donnell prove the
weaker $Z(\alpha)\subseteq[-\rho(X),\rho(X)]$, so the proposal stands to that theorem exactly as
the Conjecture of \cite{PartI} stands to Hall--Puder--Sawin.

His own remark \cite{Xu} supplies the refutation. Taking the atoms to be the individual edge-adjacency
matrices of a finite graph $G$ gives $\alpha=\mu_G$ and $\operatorname{AddProd}=\mathrm{UCT}(G)=T$,
so the conjecture specialises to $Z(\mu_G)\subseteq\operatorname{Spec}(T)$, which is
the Conjecture of \cite{PartI} at $d=1$. Hall's graph therefore refutes it, and the accompanying
$(k,m)$-characteristic-polynomial restatement falls with it through the identity
$\psi_{k,m}[k\mathcal A_{\mathrm{bd}}]=\alpha$.

\begin{proposition}\label{prop:xu}
The additive-product spectral-inclusion conjecture is false, as is its $(k,m)$ restatement.
\end{proposition}

\begin{proof}
Instantiate at the edge-atom system of the graph of Hall's theorem \cite{HallCE}. By the edge-atom
identification $\alpha=\mu_G$, so $\alpha(\sqrt5)=0$ by the factorisation in \cite{PartI}, while
$\sqrt5\notin\operatorname{Spec}(T)$ by the certificate of \cite{PartI}. The $(k,m)$ form
follows by substituting the identity above.
\end{proof}

The assembly is machine-checked in \texttt{RamaLean/XuAdditiveProduct.lean} as
\texttt{xu\_conj21\_false} and \texttt{xu\_conj24\_false}, with the edge-atom identification and
Angel--Friedman--Hoory entering as the hypotheses they are; the root and the exclusion come from
the verified \texttt{muG\_root\_sqrt5} and the certificate of Parts~B--D.

There is a dichotomy attached to that proposal which our answer sharpens. Either the conjecture
holds for all additive products, in which case the graph statement follows by the edge-atom
specialisation; or it fails somewhere, in which case real-rootedness, the free-like walk moment
identities and the spectral-radius bound are together insufficient for spectral inclusion, and
any proof of the graph statement must use a further property of edge-atom systems. The second
branch holds, and more sharply than that: the failure is already \emph{in} the edge-atom case
(\texttt{xu\_dichotomy\_second\_branch}), so no additional property of edge-atom systems can
rescue the statement, and the repair must restrict the graph class. That is what
\cite{PartI} does.


\begin{thebibliography}{9}
\bibitem{GG} C.~Godsil, I.~Gutman, \emph{On the matching polynomial of a
graph}, in: Algebraic Methods in Graph Theory, 1981.
\bibitem{HPS} C.~Hall, D.~Puder, W.~Sawin, \emph{Ramanujan coverings of
graphs}, Adv.\ Math.\ \textbf{323} (2018), 367--410.
\bibitem{CGHRS} G.~Cochran, C.~Groothuis, A.~Herring, R.~Rohatgi,
E.~Stucky, \emph{A new (combinatorial) proof of the commutativity of
matching polynomials for cycles}, arXiv:1810.05889 (2018). Preprint; no journal
reference or DOI as of August 2026.
\bibitem{WWM} J.~Wan, W.~Wang, A.~Mohammadian, \emph{On the matching
polynomials of graphs with small number of cycles}, arXiv:2103.10580 (2021).
\bibitem{SFM} Y.-M.~Song, Y.-Z.~Fan, Z.~Miao, \emph{Hypergraph coverings and
Ramanujan hypergraphs}, arXiv:2310.01771 (2023).
\bibitem{HL} O.~J.~Heilmann, E.~H.~Lieb, \emph{Theory of monomer-dimer
systems}, Comm.\ Math.\ Phys.\ \textbf{25} (1972) 190--232.
\bibitem{God81} C.~D. Godsil, \emph{Matchings and walks in graphs},
J. Graph Theory \textbf{5} (1981), 285--297.

\bibitem{God93} C.~D. Godsil, \emph{Algebraic Combinatorics}, Chapman and Hall,
New York, 1993, Theorem 6.1.1 and \S5.6.
\bibitem{MS} H.~Mizuno, I.~Sato, \emph{Characteristic polynomials of some graph
coverings}, Discrete Math.\ \textbf{142} (1995) 295--298.
\bibitem{DFMRS} C.~Dalf\'o, M.~A.~Fiol, M.~Miller, J.~Ryan, J.~\v{S}ir\'a\v{n},
\emph{An algebraic approach to lifts of digraphs}, Discrete Appl.\ Math.\
\textbf{269} (2019) 68--76.
\bibitem{MSS} A.~Marcus, D.~Spielman, N.~Srivastava, \emph{Interlacing
families I: Bipartite Ramanujan graphs of all degrees}, Ann.\ of Math.\
\textbf{182} (2015), 307--325.
\bibitem{HallCE} C.~Hall, \emph{A simple bipartite counterexample to Conjecture 10:
  spectral localization of the ordinary matching polynomial}, verification note, personal
  communication, August 2026.
\bibitem{AFH} O.~Angel, J.~Friedman, S.~Hoory, \emph{The non-backtracking spectrum of the
  universal cover of a graph}, Trans.\ Amer.\ Math.\ Soc.\ \textbf{367} (2015), 4287--4318;
  arXiv:0712.0192.
\bibitem{LMST} W.~Li, M.~Magee, M.~Sabri, D.~Thomas, \emph{Eigenvalues of maximal
  abelian covers}, arXiv:2508.17332 (2025).
\bibitem{Sp} T.~J.~Spier, \emph{Eigenvalues of universal covers and the matching
  polynomial}, arXiv:2510.05041 (2025).
\bibitem{PartI} V.~Bonfioli, \emph{Roots of $d$-matching polynomials and the spectrum of
  the universal cover: a closed form at first Betti number one, and a refuted conjecture}, 2026.
\bibitem{PartII} V.~Bonfioli, \emph{The biregular case of a spectral-inclusion
  conjecture: the weighted plane class and its obstructions}, companion paper, 2026.
\bibitem{Xu} Z.~Xu, \emph{Note on mixed characteristic polynomials}, personal communication,
  August 2026.
\bibitem{MO} S.~Mohanty, R.~O'Donnell, \emph{X-Ramanujan graphs}, arXiv:1904.03500v2 (2019).
\end{thebibliography}

\end{document}
